\section{Setup}

The corpus is synthetic by design. It consists of 97 tasks generated from 25 hand-written seed functions, where each task is a short, standard-library-only Python function into which exactly one bug has been injected by changing a single line. The bugs are drawn from a closed taxonomy of eight types, such as off-by-one errors, inverted conditions, and deleted edge-case guards. I chose synthetic tasks over organic bugs from real repositories because a synthetic corpus lets me freeze the ground truth. For every task I know the reference solution, the exact error, and the test cases that error is attributed to, so the checker's verdict is authoritative and every failure traces to the one injected bug. In an organic repository the causes of a bug are ambiguous, the correct fix is debatable, and therefore the ground truth itself is debatable. The closed taxonomy has a second benefit: because I know each task's bug class, I can later ask where the failures concentrate, not just how many there are. One honest cost of this choice is that the seed functions are common utilities the models have likely seen in training, which inflates absolute success rates but not the advisory-versus-binding difference the thesis rests on.

The checker runs a submission against the task's frozen test suite in an isolated subprocess and returns a deterministic verdict: PASSED, or FAILED with the names of the failing tests and an error excerpt. The two harnesses share everything that must be shared, including the prompt builder, the verdict renderer, and the checker itself, so the feedback text is byte-identical across arms. They differ in exactly three rules, all on the binding side: completion is computed rather than declared, a byte-identical resubmission after a failure is rejected without re-checking, and three consecutive identical failures end the episode as escalated.

The spec-starved presentation was not the starting point but a logged discovery. Initially the tasks were shown with their full descriptions, and both models solved essentially every dev task in one shot, leaving no difference for the experiment to measure. I then withheld the separate description while leaving the docstrings in the source, and still everything was one-shotted, because the docstrings alone carried enough of the specification. Only after additionally stripping all docstrings and comments from the presented code did the corpus begin to discriminate between the two modes. This presentation was adopted after the pilots but registered before any test-set episode ran.

\begin{figure}[t]
\centering
\includegraphics[width=\textwidth]{fig6_workflow.png}
\caption{The experiment workflow: boxes are states and arrows are actions. The hidden test tasks flow to the full run untouched by piloting, and the run may start only after the predictions are locked. Drawn by \texttt{writeup\_figures/plot\_diagrams.py}.}
\label{fig:workflow}
\end{figure}

The discipline around the experiment defends against a single attack: that I fit the experiment to the result. The corpus was frozen with a content hash before the runs, which proves no task was edited after the results existed. Ten tasks were reserved for pilots, and all design decisions, including the presentation changes above, were made while looking only at those ten, so the 87 test tasks stayed unopened until the design was final and cannot have been tuned toward any outcome. The predictions were committed with a git tag before the first paid API call, so the hypotheses provably precede the data. Every deviation from plan was recorded in an append-only log. Both models ran at temperature 0 with a step cap of 8, identical across arms.
